% ============================================================
%  Chương IV – Chứng chỉ số và các chuẩn
% ============================================================
\chapter{Chứng chỉ số và các chuẩn}
\label{chap:x509}

Chứng chỉ số là đối tượng dữ liệu mà PKI phát hành và
kiểm tra. Với chứng chỉ X.509, một CA ký lên dữ liệu gồm
danh tính, khóa công khai, thời hạn hiệu lực và các ràng buộc
sử dụng. Bên nhận không cần biết trước khóa công khai của chủ
thể. Họ kiểm tra chữ ký của CA, chuỗi chứng chỉ, thời hạn,
trạng thái thu hồi và tên được chứng nhận.

% ----------------------------------------------------------
\section{Chứng chỉ số X.509}
\label{sec:4.1}

X.509 xuất phát từ hệ thư mục X.500 của ITU-T. Trong Web PKI
hiện nay, cách diễn giải được dùng rộng rãi là hồ sơ Internet
X.509 do RFC~5280 định nghĩa cho chứng chỉ và CRL~\cite{RFC5280}.
Phiên bản 1 chứa các trường cơ bản. Phiên bản 2 thêm định danh
riêng cho issuer và subject. Phiên bản 3 bổ sung cơ chế
extension, nhờ đó cùng một khuôn dạng chứng chỉ có thể biểu diễn
CA, máy chủ TLS, người dùng email, thiết bị hoặc chứng chỉ ký mã.

Một chứng chỉ X.509 gồm ba phần ở mức ASN.1. Phần
\texttt{tbsCertificate} là dữ liệu \enquote{to be signed}:
CA ký đúng phần này. Phần thuật toán chữ ký cho biết thuật toán
được dùng, chẳng hạn RSA với SHA-256 hoặc ECDSA với SHA-256.
Phần giá trị chữ ký chứa chữ ký trên \texttt{tbsCertificate}.
Cấu trúc này cho phép bên kiểm tra tái tạo dữ liệu đã ký và xác
minh bằng khóa công khai của issuer.

X.509 được mã hóa theo DER trong tầng nhị phân. DER là một dạng
mã hóa xác định của ASN.1, nên cùng một cấu trúc dữ liệu có đúng
một biểu diễn byte hợp lệ. PEM là lớp bao gói văn bản: dữ liệu DER
được mã hóa Base64 và đặt giữa các dòng
\texttt{-----BEGIN CERTIFICATE-----} và
\texttt{-----END CERTIFICATE-----}. Vì vậy, DER và PEM không phải
hai loại chứng chỉ khác nhau. Chúng là hai cách biểu diễn cùng một
nội dung.

\begin{table}[ht]
    \centering
    \caption{Các trường chính trong chứng chỉ X.509 v3}
    \label{tab:x509-fields}
    \small
    \begin{tabularx}{\textwidth}{@{} >{\raggedright\arraybackslash}p{0.31\textwidth} X X @{}}
        \toprule
        \textbf{Trường} & \textbf{Ý nghĩa} & \textbf{Ghi chú} \\
        \midrule
        \texttt{version}
            & Phiên bản khuôn dạng chứng chỉ.
            & Web PKI thực tế dùng v3 để hỗ trợ extension. \\
        \texttt{serialNumber}
            & Số seri do CA gán trong phạm vi issuer.
            & Được dùng khi thu hồi qua CRL hoặc OCSP. \\
        \texttt{issuer}
            & DN của thực thể ký chứng chỉ.
            & Thường là CA trung gian hoặc root CA. \\
        \texttt{validity}
            & Cặp thời điểm \texttt{notBefore} và \texttt{notAfter}.
            & Chứng chỉ chỉ hợp lệ trong khoảng này. \\
        \texttt{subject}
            & DN của chủ thể được cấp chứng chỉ.
            & Có thể rỗng nếu danh tính nằm trong SAN. \\
        \texttt{subjectPublicKeyInfo}
            & Thuật toán khóa công khai và giá trị khóa.
            & Ví dụ RSA, ECDSA hoặc EdDSA. \\
        \texttt{extensions}
            & Các ràng buộc và thuộc tính mở rộng.
            & Là cơ chế trung tâm của X.509 v3. \\
        \bottomrule
    \end{tabularx}
\end{table}

\section{Các trường X.509 v3}
\label{sec:4.2}

Trường \texttt{version} xác định cách diễn giải chứng chỉ. Với
X.509 v3, trường \texttt{extensions} trở thành bắt buộc về mặt
thực tiễn vì phần lớn kiểm tra an toàn nằm ở đây. Trường
\texttt{serialNumber} phải đủ phân biệt trong phạm vi issuer để
có thể định danh chứng chỉ khi thu hồi. Trường \texttt{issuer}
cho biết tên của CA ký chứng chỉ; trường \texttt{subject} cho biết
tên của thực thể nhận chứng chỉ.

Trường \texttt{validity} gồm \texttt{notBefore} và
\texttt{notAfter}. Bên kiểm tra phải từ chối chứng chỉ nếu thời
điểm hiện tại nằm ngoài khoảng này. Trường
\texttt{subjectPublicKeyInfo} chứa khóa công khai của subject cùng
định danh thuật toán. Với chứng chỉ máy chủ, đây là khóa được dùng
để xác thực chữ ký hoặc tham gia cơ chế trao đổi khóa tùy giao
thức phía trên. Chi tiết bắt tay TLS không được trình bày ở đây,
vì thuộc phạm vi Chương~\ref{chap:deploy}.

\section{Distinguished Name}
\label{sec:4.3}

Distinguished Name (DN) là tên phân cấp theo mô hình X.500.
Trong chứng chỉ, DN thường xuất hiện ở \texttt{issuer} và
\texttt{subject}. Các thuộc tính thường gặp gồm
\texttt{C} cho quốc gia, \texttt{ST} cho bang hoặc tỉnh,
\texttt{L} cho địa phương, \texttt{O} cho tổ chức,
\texttt{OU} cho đơn vị tổ chức, \texttt{CN} cho common name và
\texttt{emailAddress} cho địa chỉ thư điện tử. DN hữu ích trong
môi trường tổ chức vì nó mô tả subject bằng các thuộc tính có cấu
trúc.

Trong HTTPS, \texttt{CN} không còn là nơi chính để kiểm tra tên
miền. RFC~6125 yêu cầu ứng dụng ưu tiên subject alternative name
và không quay lại \texttt{CN} khi SAN thuộc loại tên phù hợp đã có
mặt~\cite{RFC6125}. Điều này tách rõ tên hiển thị trong DN khỏi
tập tên mà chứng chỉ thực sự chứng nhận. Với chứng chỉ máy chủ,
SAN mới là nơi chứa DNS name hoặc IP address được so khớp với tên
mà client truy cập.

\section{Extension quan trọng}
\label{sec:4.4}

Extension làm cho X.509 v3 phù hợp với nhiều ngữ cảnh sử dụng.
Mỗi extension có định danh OID, giá trị và cờ critical. Nếu một
extension critical không được hiểu, bên kiểm tra phải từ chối
chứng chỉ theo RFC~5280~\cite{RFC5280}. Cơ chế này cho phép CA áp
đặt điều kiện an toàn mà client không được bỏ qua.

\begin{table}[ht]
    \centering
    \caption{Một số extension thường gặp trong X.509 v3}
    \label{tab:x509-extensions}
    \small
    \begin{tabularx}{\textwidth}{@{} >{\raggedright\arraybackslash}p{0.31\textwidth} X X @{}}
        \toprule
        \textbf{Extension} & \textbf{Nội dung} & \textbf{Vai trò} \\
        \midrule
        \texttt{BasicConstraints}
            & Cờ \texttt{cA} và tùy chọn \texttt{pathLenConstraint}.
            & Phân biệt CA với end-entity certificate. \\
        \texttt{KeyUsage}
            & Bit như \texttt{digitalSignature}, \texttt{keyEncipherment},
              \texttt{keyCertSign}.
            & Giới hạn thao tác mật mã được phép. \\
        \texttt{ExtendedKeyUsage}
            & OID mục đích như server authentication, client authentication
              hoặc code signing.
            & Thu hẹp ngữ cảnh sử dụng của khóa. \\
        \texttt{SubjectAltName}
            & DNS name, IP address, email hoặc URI của subject.
            & Cơ sở kiểm tra tên miền trong Web PKI. \\
        \texttt{SubjectKeyIdentifier}
            & Định danh suy ra từ khóa công khai của subject.
            & Hỗ trợ tìm và nối chuỗi chứng chỉ. \\
        \texttt{AuthorityKeyIdentifier}
            & Định danh khóa của issuer.
            & Ghép chứng chỉ với CA đã ký nó. \\
        \texttt{CRLDistributionPoints}
            & Vị trí tải CRL.
            & Cho phép client kiểm tra thu hồi bằng CRL. \\
        \texttt{AuthorityInfoAccess}
            & Vị trí OCSP responder hoặc chứng chỉ issuer.
            & Hỗ trợ kiểm tra trạng thái và dựng chuỗi. \\
        \bottomrule
    \end{tabularx}
\end{table}

\texttt{BasicConstraints} là extension bắt buộc về mặt an toàn với
CA certificate. Nếu \texttt{cA} bằng \texttt{TRUE}, chứng chỉ có
thể được dùng để ký chứng chỉ khác, với điều kiện các trường khác
như \texttt{KeyUsage} cũng cho phép. \texttt{pathLenConstraint}
giới hạn số CA trung gian có thể xuất hiện bên dưới CA đó trong
chuỗi chứng chỉ.

\texttt{KeyUsage} mô tả thao tác mật mã cấp thấp. Ví dụ,
\texttt{digitalSignature} cho phép dùng khóa để kiểm tra chữ ký số,
trong khi \texttt{keyCertSign} chỉ phù hợp với CA. Ngược lại,
\texttt{ExtendedKeyUsage} mô tả mục đích ứng dụng. Một chứng chỉ
có thể chứa khóa hợp lệ về mặt thuật toán nhưng vẫn bị từ chối nếu
EKU không chứa mục đích mà ứng dụng cần.

\texttt{SubjectAltName} chứa các tên được chứng nhận cho subject.
Với chứng chỉ website, SAN thường chứa một hoặc nhiều DNS name.
\texttt{SubjectKeyIdentifier} và \texttt{AuthorityKeyIdentifier}
không tự tạo ra niềm tin, nhưng giúp phần mềm chọn đúng issuer khi
có nhiều chứng chỉ CA cùng tên hoặc cùng tổ chức. Hai extension
\texttt{CRLDistributionPoints} và \texttt{AuthorityInfoAccess}
liên kết chứng chỉ với hạ tầng kiểm tra thu hồi qua CRL và
OCSP~\cite{RFC6960}.

\section{Certificate Signing Request}
\label{sec:4.5}

Certificate Signing Request (CSR) là yêu cầu ký chứng chỉ do chủ
thể gửi cho CA. Khuôn dạng CSR phổ biến là PKCS~\#10, được RFC
2986 mô tả~\cite{RFC2986}. CSR chứa thông tin subject, khóa công
khai, một tập thuộc tính tùy chọn và chữ ký do chính khóa bí mật
ứng với khóa công khai trong CSR tạo ra. Chữ ký này chứng minh
rằng người gửi CSR có quyền kiểm soát khóa bí mật, nhưng nó không
chứng minh danh tính pháp lý hoặc quyền sở hữu tên miền. Phần đó
thuộc trách nhiệm xác minh của RA hoặc CA.

Luồng cấp chứng chỉ có bốn bước. Thứ nhất, subject tạo cặp khóa và
bảo vệ khóa bí mật cục bộ. Thứ hai, subject tạo CSR chứa khóa công
khai và thông tin định danh mong muốn. Thứ ba, CA kiểm tra quyền
kiểm soát tên miền, thông tin tổ chức hoặc tiêu chí tương ứng với
loại chứng chỉ. Thứ tư, CA ký dữ liệu X.509 và trả lại chứng chỉ.
Khóa bí mật không cần rời khỏi hệ thống của subject trong luồng
thông thường.

\section{Phân loại chứng chỉ}
\label{sec:4.6}

CA/Browser Forum Baseline Requirements định nghĩa yêu cầu tối
thiểu cho chứng chỉ máy chủ TLS được tin cậy công khai, bao gồm
quy trình xác minh tên miền, thông tin subject và vận hành CA
\cite{CABForum}. Trong thực tế, chứng chỉ TLS thường được phân
loại theo mức xác minh danh tính: domain validation (DV),
organization validation (OV) và extended validation (EV). DV xác
minh quyền kiểm soát tên miền. OV thêm kiểm tra tổ chức. EV áp
dụng bộ yêu cầu nhận diện tổ chức chặt hơn theo tài liệu EV riêng
của CA/Browser Forum~\cite{CABForumEV}.

\begin{table}[ht]
    \centering
    \caption{Phân loại chứng chỉ theo mục đích sử dụng}
    \label{tab:cert-types}
    \small
    \begin{tabularx}{\textwidth}{@{} >{\raggedright\arraybackslash}p{0.31\textwidth} X X @{}}
        \toprule
        \textbf{Loại} & \textbf{Đặc điểm} & \textbf{Cách dùng} \\
        \midrule
        DV
            & Xác minh quyền kiểm soát tên miền.
            & Website công khai, tự động hóa cấp phát. \\
        OV
            & Xác minh thêm thông tin tổ chức.
            & Dịch vụ cần hiện thông tin pháp nhân rõ hơn. \\
        EV
            & Xác minh tổ chức theo yêu cầu EV.
            & Môi trường cần quy trình nhận diện chặt hơn. \\
        Wildcard
            & Tên dạng \texttt{*.example.com}.
            & Nhiều subdomain cùng một cấp. \\
        Multi-domain/SAN
            & Nhiều tên trong \texttt{SubjectAltName}.
            & Gom nhiều DNS name vào một chứng chỉ. \\
        Code signing
            & EKU cho ký mã hoặc chính sách tương ứng.
            & Ký phần mềm, script hoặc gói cài đặt. \\
        S/MIME
            & Gắn khóa công khai với địa chỉ email.
            & Ký và mã hóa thư điện tử~\cite{RFC8551}. \\
        Client certificate
            & Subject đại diện người dùng, thiết bị hoặc dịch vụ.
            & Xác thực client trong VPN, mTLS hoặc hệ nội bộ. \\
        \bottomrule
    \end{tabularx}
\end{table}

Wildcard certificate chỉ bao phủ một cấp nhãn theo cách diễn giải
thông thường của Web PKI. Ví dụ, \texttt{*.example.com} có thể
khớp \texttt{www.example.com}, nhưng không khớp
\texttt{a.b.example.com}. Multi-domain certificate dùng SAN để
chứa nhiều tên cụ thể. Cách này phù hợp khi một dịch vụ có nhiều
tên miền hoặc khi cần hợp nhất chứng chỉ cho các endpoint liên
quan.

\section{Các chuẩn PKCS}
\label{sec:4.7}

PKCS là họ chuẩn ban đầu do RSA Laboratories công bố cho các đối
tượng mật mã trong triển khai thực tế. Trong phạm vi chứng chỉ số,
ba chuẩn cần quan tâm là PKCS~\#7/CMS, PKCS~\#10 và PKCS~\#12.
Báo cáo không liệt kê toàn bộ họ PKCS vì nhiều chuẩn đã lỗi thời
hoặc nằm ngoài trọng tâm PKI.

PKCS~\#7 định nghĩa cú pháp cho dữ liệu được ký hoặc mã hóa. Hậu
thân chuẩn hóa trong IETF là Cryptographic Message Syntax (CMS),
được RFC~5652 mô tả~\cite{RFC5652}. CMS là nền tảng cho nhiều
định dạng chữ ký và thông điệp bảo mật, trong đó có S/MIME.
PKCS~\#10 là định dạng CSR như đã trình bày ở
Mục~\ref{sec:4.5}. PKCS~\#12 định nghĩa định dạng đóng gói khóa
bí mật, chứng chỉ và chuỗi chứng chỉ vào một tệp, thường có phần
bảo vệ bằng mật khẩu; RFC~7292 mô tả định dạng này~\cite{RFC7292}.
Tệp \texttt{.p12} hoặc \texttt{.pfx} thường được dùng khi xuất
nhập chứng chỉ người dùng hoặc chứng chỉ máy chủ giữa các hệ thống.

\section{RFC và chuẩn quan trọng}
\label{sec:4.8}

RFC~5280 là chuẩn nền tảng cho hồ sơ Internet X.509, bao gồm cấu
trúc chứng chỉ, extension, thuật toán dựng đường chứng chỉ và CRL
\cite{RFC5280}. RFC~2986 định nghĩa PKCS~\#10 cho CSR
\cite{RFC2986}. RFC~5652 định nghĩa CMS, là khuôn dạng chung cho
dữ liệu ký số và mã hóa ở tầng thông điệp~\cite{RFC5652}. RFC~7292
định nghĩa PKCS~\#12 cho gói khóa và chứng chỉ~\cite{RFC7292}.

RFC~6960 định nghĩa OCSP, giao thức hỏi trạng thái thu hồi cho một
chứng chỉ cụ thể~\cite{RFC6960}. RFC~6962 định nghĩa Certificate
Transparency phiên bản thử nghiệm, trong đó chứng chỉ TLS được ghi
vào log công khai dựa trên cây Merkle~\cite{RFC6962}. Các triển
khai mới hơn có thể tham chiếu RFC~9162 cho Certificate
Transparency phiên bản 2.0~\cite{RFC9162}. CA/Browser Forum
Baseline Requirements không thay thế RFC, nhưng đặt yêu cầu vận
hành cho CA muốn được tin cậy trong hệ sinh thái trình duyệt
\cite{CABForum}. Vì vậy, Web PKI là kết quả của cả chuẩn kỹ thuật
và quy tắc quản trị CA.
